\section{The congruence complexes}
\label{sec:construction}

We now define the complexes appearing in \cref{thm:main}. Fix
\[
 n=2d,\qquad N=n+1=2d+1,
\]
and choose a sufficiently large odd prime $q$, depending only on $N$.
The precise largeness requirements will be stated where they are used.
Set $k=\mathbb F_q$, $R=k[t]$, and
\[
 \Gamma_a(m)=\ker\bigl(\SL_a(R)\longrightarrow\SL_a(R/(t^m))\bigr).
\]
Let $U\le\SL_N(R)$ be the inverse image, under evaluation at $t=0$, of
the upper unitriangular subgroup of $\SL_N(k)$. Put
\[
 S=\{0,1,\ldots,N-1\},\qquad
 x_i=e_{i,i+1}(1)\ (1\le i<N),\qquad x_0=e_{N,1}(t),
\]
where $e_{ab}(f)=I+fE_{ab}$, and define
\begin{equation}\label{eq:Hi}
 H_i=\langle x_j:j\in S\setminus\{i\}\rangle.
\end{equation}

For a group $G$ and subgroups $(K_i)_{i\in S}$, the \emph{coset complex}
$\CC(G;(K_i))$ has vertices the tagged left cosets $gK_i$; a collection
of vertices is a face when the corresponding cosets have nonempty common
intersection. Its chambers are the sets $\{gK_i:i\in S\}$, so it is pure
of dimension $|S|-1$. For $m\ge2$, set
\begin{equation}\label{eq:construction}
 G_m=U/\Gamma_N(m),\qquad
 T_m=\CC(G_m;(\overline H_i)_{i\in S}),\qquad
 X_m=\skel_d T_m.
\end{equation}
The bar denotes reduction modulo $t^m$. Here $\Gamma_N(m)\subseteq U$,
and $G_m$ is the finite image of $U$ in $\SL_N(R/(t^m))$.

